\documentclass[11pt]{article}
\usepackage{amsmath,amsthm,amssymb}
\usepackage[margin=1in]{geometry}

\newtheorem{theorem}{Theorem}
\newtheorem{lemma}{Lemma}
\newtheorem{proposition}{Proposition}
\newtheorem{corollary}{Corollary}
\newtheorem{conjecture}{Conjecture}
\theoremstyle{definition}
\newtheorem{definition}{Definition}
\theoremstyle{remark}
\newtheorem{remark}{Remark}

\DeclareMathOperator{\Des}{Des}
\DeclareMathOperator{\SYT}{SYT}
\DeclareMathOperator{\ord}{ord}
\DeclareMathOperator{\supp}{supp}
\DeclareMathOperator{\Real}{Re}
\DeclareMathOperator{\Imag}{Im}
\newcommand{\ZZ}{\mathbb{Z}}
\newcommand{\CC}{\mathbb{C}}
\newcommand{\la}{\lambda}
\newcommand{\inner}[1]{\left\langle #1 \right\rangle}
\newcommand{\psib}{\overline{\psi}}

\title{The $d=4$ fiber of the graded order law:\\
a reduction to Schur support, and the obstruction}
\author{Clio}
\date{2026-06-03}

\begin{document}
\maketitle

\begin{abstract}
Let $G_\la(q)=\sum_{T\in\SYT(\la)}q^{s(T)}$ be the graded branch--exponent
generating function, where $s(T)=\sum_{i\in\Des(T)}w_i$ with the parity weights
$w_i=2i-1$ if $n-i$ is odd and $w_i=0$ otherwise. At the cube root of unity the
fiber value $G_\la(\zeta_3)$ vanishes exactly on the two $3$-cores $(3,1),(2,1,1)$,
by a Schur-positivity argument [Clio 2026-06-02]. At the \emph{fourth} root of
unity $\zeta_4=i$ the value is genuinely complex and that argument cannot
transplant. We prove instead a structural \emph{reduction}: because the active
descent set is an arithmetic progression of common difference $2$, at $\zeta_4$
\emph{all} of its factors in the master identity coincide, so that for
$n=2m$ even
\[
   G_\la(i)=\inner{s_\la,\ \psi^{\,m}},\qquad \psi:=h_2+i\,e_2=s_2+i\,s_{11},
\]
a single pure power, with the odd-$n$ analogue $G_\la(i)=\inner{s_\la,h_1\psib^{\,m}}$.
Equivalently $G_\la(i)=P_\la(i)$ for the nonnegative-integer, conjugation-reciprocal
polynomial $P_\la(t)=\inner{s_\la,(h_2+te_2)^m}$, so $G_\la(i)=0\iff(t^2+1)\mid
P_\la(t)\iff\la$ lies outside the (complex) Schur support of $\psi^m$. We then prove
two complete results. \emph{(i) A parity criterion:} in $\ZZ[i]$ modulo the prime
$1+i$, $G_\la(i)\equiv i^m f^\la$, so \emph{$f^\la$ odd implies $G_\la(i)\neq0$}, with
the refinement $\Real G_\la(i)\equiv\tfrac12(f^\la+\chi^\la(2^m))$, $\Imag
G_\la(i)\equiv\tfrac12(f^\la-\chi^\la(2^m))\pmod2$ and explicit mod-$4$ necessary
conditions for a vanisher. \emph{(ii) Exact values} from a weight-$\{1,i,1+i\}$
$2$-strip recursion: $G_{(n)}=1$, $G_{(1^n)}=i^m$, $G_{(2M-1,1)}=(M-1)+Mi$, and
\[
   \boxed{\,G_{(2m-2,2)}(i)=m(m-2)\,i\,},
\]
the last vanishing \emph{exactly} at $m=2$, so within two-row shapes $(2,2)$ is
provably the unique vanisher. The general case reduces to full support of $\psi^m$
(for $m\ge3$; for $m=2$ the support omits exactly $(2,2)$, a one-line expansion). We
isolate why this is hard -- $\psi^m$ has genuinely complex coefficients, the Pieri
recursion admits cancellation, and the Schur-positivity that closed $d=3$ is absent --
and give two precise routes (a $(1+i)$-adic monovariant; a Hermitian bound from
$\psi^m\overline{\psi^m}=\Xi^m$), each with its remaining obstruction.
\end{abstract}

\section{Setup}

Throughout $\la\vdash n$, and for $T\in\SYT(\la)$,
$\Des(T)=\{i\in[n-1]:i{+}1\text{ lies in a strictly lower row than }i\}$. Fix
\[
  A=\{k\in[n-1]:n-k\text{ odd}\},\qquad w_k=(2k-1)\,[k\in A],\qquad
  s(T)=\sum_{i\in\Des(T)}w_i,
\]
and set $m=|A|=\lfloor n/2\rfloor$, $e=n\bmod2$. We work in the ring $\Lambda$ of
symmetric functions with the Hall inner product, for which
$\inner{s_\la,h_\mu}=K_{\la\mu}$ and $\inner{s_\la,p_\mu}=\chi^\la(\mu)$. We use
$h_2=\tfrac12(p_1^2+p_2)$, $e_2=\tfrac12(p_1^2-p_2)$, so $h_2+e_2=p_1^2$ and
$h_2-e_2=p_2$. For $f\in\Lambda$ write $\supp_\CC(f)=\{\mu:\inner{s_\mu,f}\neq0\}$.

We recall the master identity, proved (for any $\zeta$) in the $\zeta_3$ note.

\begin{proposition}[Master identity]\label{prop:master}
For any $\zeta\in\CC$,
$\displaystyle
  G_\la(\zeta)=\Big\langle s_\la,\ h_1^{\,e}\prod_{k\in A}\big(h_2+e_2\,\zeta^{\,2k-1}\big)\Big\rangle .
$
\end{proposition}

\section{The reduction: a single pure power at $\zeta_4$}

This is the structural phenomenon that distinguishes $d=4$ from general $d$.

\begin{lemma}[Coincidence of factors]\label{lem:coincide}
At $\zeta=i$, for every $k\in A$ the factor $h_2+e_2 i^{2k-1}$ equals
$h_2+i\,e_2$ if $n$ is even, and $h_2-i\,e_2$ if $n$ is odd. In particular all
$|A|=m$ factors are equal.
\end{lemma}

\begin{proof}
Since $2k-1$ is odd, $i^{2k-1}=i$ if $2k-1\equiv1\pmod4$, i.e.\ $k$ odd, and
$i^{2k-1}=i^3=-i$ if $k$ even. The active set $A$ is an arithmetic progression of
common difference $2$: $A=\{1,3,\dots,n-1\}$ when $n$ is even (all $k$ odd) and
$A=\{2,4,\dots,n-1\}$ when $n$ is odd (all $k$ even). Hence within $A$ the parity
of $k$ is constant, and $i^{2k-1}$ takes a single value: $i$ for $n$ even, $-i$
for $n$ odd.
\end{proof}

\begin{theorem}[Reduction]\label{thm:reduction}
Write $\psi=h_2+i\,e_2$ and $\psib=h_2-i\,e_2$. Then
\[
  n=2m\ \text{even:}\quad G_\la(i)=\inner{s_\la,\ \psi^{\,m}},
  \qquad\qquad
  n=2m+1\ \text{odd:}\quad G_\la(i)=\inner{s_\la,\ h_1\,\psib^{\,m}}.
\]
Consequently, for $n$ even, $G_\la(i)=0$ if and only if
$\la\notin\supp_\CC(\psi^m)$.
\end{theorem}

\begin{proof}
Immediate from Proposition~\ref{prop:master} and Lemma~\ref{lem:coincide}: the
product $\prod_{k\in A}(h_2+e_2 i^{2k-1})$ collapses to $\psi^{m}$ (resp.\
$\psib^{m}$), and $h_1^{e}=1$ (resp.\ $h_1$). The support statement is the
definition of $\inner{s_\la,\psi^m}\neq0$.
\end{proof}

\begin{remark}
This is exactly the point where the $\zeta_3$ proof's strategy breaks. At
$\zeta_3$ the factors split into a conjugate pair $\phi,\bar\phi$ whose product
$\Theta=\phi\bar\phi=s_4+2s_{22}+s_{1^4}$ is \emph{Schur-positive}, turning the
fiber value into a manifestly nonnegative integer. At $\zeta_4$ with $n$ even
\emph{all} factors coincide as $\psi$; there is no conjugate partner inside the
product, so $\psi^m$ has genuinely complex coefficients and no positivity is
available. The value $G_\la(i)$ is a Gaussian integer that can a priori vanish by
cancellation.
\end{remark}

\noindent We henceforth treat the even case $n=2m$ and write
$P_\la:=\inner{s_\la,\psi^m}=G_\la(i)\in\ZZ[i]$. The odd case is recovered from
$G_\la(i)=\inner{s_\la,h_1\psib^m}=\sum_{\mu}\,\overline{P_\mu}$, the sum over the
even partitions $\mu$ obtained by deleting one corner box of $\la$.

\section{Three faces of $P_\la$}

\subsection{The interpolating polynomial and the skewing recursion}
Let $P_\la(t):=\inner{s_\la,(h_2+t\,e_2)^m}\in\ZZ_{\ge0}[t]$, so $P_\la=P_\la(i)$.
Because $\inner{s_\la,(h_2+te_2)^m}$ is a product of the Schur-positive functions
$h_2,e_2$, its coefficients are nonnegative integers. The skewing operators
$h_2^{\perp},e_2^{\perp}$ remove horizontal, resp.\ vertical, $2$-strips, giving
the recursion
\begin{equation}\label{eq:rec}
  P_\la(t)=\sum_{\la/\mu\ \text{horiz.\ }2\text{-strip}}P_\mu(t)
          \;+\;t\!\!\sum_{\la/\mu\ \text{vert.\ }2\text{-strip}}\!\!P_\mu(t),
  \qquad P_\varnothing=1.
\end{equation}
Three evaluations are pinned down: $P_\la(1)=\inner{s_\la,p_1^{2m}}=f^\la$;
$P_\la(-1)=\inner{s_\la,p_2^{m}}=\chi^\la(2^m)=\pm\binom{m}{|\la^0|}f^{\la^0}f^{\la^1}$
(the signed number of standard domino tableaux, nonzero iff $\la$ has empty
$2$-core); and $P_\la(0)=K_{\la,(2^m)}$. Conjugation acts by coefficient
reversal,
\begin{equation}\label{eq:conj}
  P_{\la'}(t)=t^{m}P_\la(1/t),\qquad\text{hence}\qquad
  P_{\la'}(i)=i^{m}\,\overline{P_\la(i)},
\end{equation}
using $\omega(h_2+te_2)=t\,(h_2+t^{-1}e_2)$. In particular $G_\la(i)=0\iff
G_{\la'}(i)=0$, and the self-conjugate $(2,2)$ is consistent with this.

\subsection{Hermite determinant and Gaussian integral}
Setting $p_{\ge3}=0$ is the ring homomorphism $\Lambda\to\CC[A,B]$,
$p_1\mapsto A,\ p_2\mapsto B$. Under it
$h_r\mapsto[z^r]\exp\!\big(Az+\tfrac{B}{2}z^2\big)$, an Appell/Hermite polynomial,
so by Jacobi--Trudi $s_\la(p_1{=}A,p_2{=}B,0)=\det\big(h_{\la_i-i+j}(A,B,0)\big)$
is a Hermite determinant. A Hubbard--Stratonovich linearisation of $p_1^2$ then
yields the integral representation
\begin{equation}\label{eq:gauss}
  P_\la(t)=\frac1{\sqrt\pi}\int_{-\infty}^{\infty} e^{-s^2}\,
           s_\la\big(p_1{=}2s,\ p_2{=}2t,\ p_{\ge3}{=}0\big)\,ds ,
\end{equation}
equivalently $G_\la(i)=\mathbb E_{Y\sim N(0,2)}\big[s_\la(p_1{=}Y,p_2{=}-2i,0)\big]$
up to the unit $\beta^m i^m$, $\beta=\tfrac{1-i}2$. (Verified numerically for all
$\la$ with $n\le10$.)

\subsection{A Parseval identity}
By orthonormality of Schur functions,
\begin{equation}\label{eq:parseval}
  \sum_{\la\vdash2m}\big|G_\la(i)\big|^2
   =\inner{\psi^m,\psib^{\,m}}
   =2^{-m}\sum_{a=0}^{m}\binom{m}{a}^{2}2^{a}\,a!\,(2m-2a)!\ >0,
\end{equation}
since $\inner{p_\mu,p_\nu}=z_\mu\delta_{\mu\nu}$ and $\alpha\beta=\tfrac12$ for
$\psi=\alpha p_1^2+\beta p_2$ ($\alpha=\tfrac{1+i}2$). (Verified for $m\le5$:
the totals are $2,12,168,4512,198720$.) This shows the fiber values are not all
zero but does not by itself isolate $(2,2)$.

\section{What is proved, and the boundary}

\begin{proposition}[Extreme shapes]\label{prop:extreme}
$P_{(2m)}(i)=1$ and $P_{(1^{2m})}(i)=i^{m}$; both are nonzero.
\end{proposition}
\begin{proof}
A one-row shape admits only horizontal $2$-strip removals, and there is a unique
chain $(2m)\to(2m-2)\to\cdots\to\varnothing$; by \eqref{eq:rec} its weight is
$t^{0}$, so $P_{(2m)}(t)=1$. Dually $P_{(1^{2m})}(t)=t^{m}$ by \eqref{eq:conj}.
\end{proof}

\begin{proposition}[Kostka expansion and a Gaussian-prime obstruction]\label{prop:kostka}
Since $e_2=h_1^2-h_2$, one has $\psi=i\,h_1^2+(1-i)h_2$, hence
\[
  \psi^m=\sum_{k=0}^m\binom mk i^{\,m-k}(1-i)^k\,h_{(2^k1^{2m-2k})},
  \qquad\text{so}\qquad
  G_\la(i)=i^m\sum_{k=0}^m\binom mk(-1-i)^k\,K_{\la,(2^k1^{2m-2k})} .
\]
Let $(1+i)$ be the Gaussian prime above $2$. Reducing modulo $(1+i)$ kills every
$k\ge1$ term, leaving $G_\la(i)\equiv i^m K_{\la,(1^{2m})}=i^m f^\la\pmod{(1+i)}$.
Therefore
\[
  f^\la\ \text{odd}\ \Longrightarrow\ G_\la(i)\neq0 .
\]
In particular every vanisher has $f^\la$ even; consistent with $f^{(2,2)}=2$.
(Both displays verified for all $\la$ with $n\le12$, $0$ exceptions; the
sufficient condition has no exception through $n=14$.)
\end{proposition}
\begin{proof}
The expansion is the binomial theorem applied to $\psi=i\,h_1^2+(1-i)h_2$ together
with $h_1^{2(m-k)}h_2^k=h_{(2^k1^{2m-2k})}$; pairing with $s_\la$ uses
$\inner{s_\la,h_\mu}=K_{\la\mu}$ and $i^{m-k}(1-i)^k=i^m\big((1-i)/i\big)^k
=i^m(-1-i)^k$. Since $(-1-i)=-(1+i)$, each $k\ge1$ summand lies in
$(1+i)\ZZ[i]$, so $G_\la(i)\equiv i^m f^\la\pmod{(1+i)}$. As $(1+i)\mid f^\la$ in
$\ZZ[i]$ iff $2\mid f^\la$, an odd $f^\la$ forces $G_\la(i)\not\equiv0$, hence
$G_\la(i)\neq0$.
\end{proof}

\begin{remark}
The obstruction is sharp at the level of one prime: it cannot see vanishers,
which all have $f^\la$ even. One might hope to iterate $(1+i)$-adically, but the
even- and odd-$k$ partial sums do \emph{not} separately vanish at a vanisher
(for $(2,2)$ they are $2+2i$ and $-2-2i$, which cancel), so a single congruence
does not suffice; controlling $v_{(1+i)}(G_\la(i))$ in general appears to require
the $2$-adic valuations of the Kostka numbers $K_{\la,(2^k1^{2m-2k})}$ (the
$2$-core tower).
\end{remark}

The single congruence of Proposition~\ref{prop:kostka} refines, at no extra cost,
into a statement about the real and imaginary parts separately, which sharpens the
necessary conditions for a vanisher.

\begin{proposition}[Refined parity congruence]\label{prop:refined}
Write $G_\la(i)=a+bi$ and $\chi:=\chi^\la(2^m)$. Then
\[
   a=\Real G_\la(i)\equiv\tfrac12\big(f^\la+\chi\big),\qquad
   b=\Imag G_\la(i)\equiv\tfrac12\big(f^\la-\chi\big)\pmod 2 .
\]
\end{proposition}
\begin{proof}
By \eqref{eq:rec}, $G_\la(i)=\sum_{v\ge0}a_v\,i^{v}$ with $a_v\in\ZZ_{\ge0}$.
Grouping by $v\bmod4$ and writing $n_r=\sum_{v\equiv r\,(4)}a_v$, we have
$a=n_0-n_2\equiv n_0+n_2\pmod2$ and $b=n_1-n_3\equiv n_1+n_3\pmod2$. Now
$n_0+n_2=\sum_{v\text{ even}}a_v=:N_{\mathrm e}$ and $n_1+n_3=\sum_{v\text{
odd}}a_v=:N_{\mathrm o}$ satisfy $N_{\mathrm e}+N_{\mathrm o}=P_\la(1)=f^\la$ and
$N_{\mathrm e}-N_{\mathrm o}=P_\la(-1)=\chi$, so $N_{\mathrm e}=\tfrac12(f^\la+\chi)$,
$N_{\mathrm o}=\tfrac12(f^\la-\chi)$.
\end{proof}

\begin{corollary}[Mod-$4$ necessary conditions]\label{cor:nec}
If $G_\la(i)=0$ then $f^\la$ and $\chi^\la(2^m)$ are even and $f^\la\pm\chi^\la(2^m)
\equiv0\pmod4$. In particular, if $\la$ has nonempty $2$-core then
$\chi^\la(2^m)=0$ and one needs $4\mid f^\la$; if $\la$ has empty $2$-core then, by
the Stanton--White formula \emph{[Clio 2026-06-02, Thm.~1]},
$\chi^\la(2^m)=(-1)^{n(\la)}f^{\la^{0}}f^{\la^{1}}\binom{m}{|\la^{0}|}$, and one
needs $f^\la\pm\chi^\la(2^m)\equiv0\pmod4$.
\end{corollary}
\begin{proof}
$G_\la(i)=0$ forces $a=b=0$, so both $\tfrac12(f^\la\pm\chi)$ are even by
Proposition~\ref{prop:refined}; hence $f^\la\pm\chi\equiv0\pmod4$, forcing
$f^\la,\chi$ even. The class $(2^m)$ has no fixed points, so $\chi^\la(2^m)\neq0$
only when $\la$ is tileable by $m$ dominoes, i.e.\ has empty $2$-core; the
empty-core value is the cited product ($c_0=0$).
\end{proof}

These congruences exclude infinitely many shapes (all odd-$f^\la$, and more), but
infinitely many shapes still satisfy them, so the parity criterion alone cannot
terminate.

\begin{proposition}[The $m=2$ support, exactly]\label{prop:m2}
$\psi^2=s_{4}+(1+2i)s_{31}+(2i-1)s_{211}-s_{1^4}$. Hence
$\supp_\CC(\psi^2)=\{(4),(3,1),(2,1,1),(1^4)\}=\{\la\vdash4\}\setminus\{(2,2)\}$,
and $(2,2)$ is the unique partition of $4$ with $G_\la(i)=0$.
\end{proposition}
\begin{proof}
With $h_2=s_2,\ e_2=s_{11}$ and Pieri, $s_2^2=s_4+s_{31}+s_{22}$,
$s_{11}^2=s_{22}+s_{211}+s_{1^4}$, $s_2 s_{11}=s_{31}+s_{211}$. Therefore
$\psi^2=(s_2+is_{11})^2=s_2^2+2is_2 s_{11}-s_{11}^2
=s_4+(1+2i)s_{31}+(1-1)s_{22}+(2i-1)s_{211}-s_{1^4}$.
The coefficient of $s_{22}$ is $1-1=0$: the real parts coming from $s_2^2$ and
$-s_{11}^2$ cancel, and the imaginary cross term $2is_2 s_{11}$ contributes
nothing to $s_{22}$. Equivalently $G_{(2,2)}(i)=\chi^{(2,2)}(1^4)\alpha^2+
\chi^{(2,2)}(2,2)\beta^2=2(\alpha^2+\beta^2)=2\cdot0=0$ since
$\alpha^2+\beta^2=(\alpha+\beta)^2-2\alpha\beta=1-1=0$.
\end{proof}

\subsection{Exact values on hooks and two rows}

The skewing recursion \eqref{eq:rec} stays inside the hook and two-row families and
solves there in closed form. We use the weights of a $2$-strip removal $\la/\mu$
read off from \eqref{eq:rec}: $1$ if $\la/\mu$ is a connected horizontal domino, $i$
if connected vertical, and $1+i$ if $\la/\mu$ is a disconnected pair (distinct row
\emph{and} column), the last because such a pair is simultaneously a horizontal and a
vertical $2$-strip.

\begin{lemma}[Near-rectangular hook]\label{lem:21}
For $M\ge1$ \emph{(}$n=2M$\emph{)}, $\,G_{(2M-1,1)}(i)=(M-1)+M\,i$, which is never
zero.
\end{lemma}
\begin{proof}
From $(a,1)$ with $a=2M-1$ the removable $2$-strips are: the arm pair
$(1,a),(1,a-1)$ (horizontal, $\to(a-2,1)$, weight $1$, present iff $a\ge3$), and the
disconnected pair $(1,a),(2,1)$ (distinct row and column as $a\ge2$, $\to(a-1)$,
weight $1+i$). With $G_{(a-1)}=1$ (Prop.~\ref{prop:extreme}) this gives, for $M\ge2$,
$G_{(2M-1,1)}=G_{(2M-3,1)}+(1+i)$. The base $G_{(1,1)}=i$ (a single vertical domino)
solves the recursion to $i+(M-1)(1+i)=(M-1)+Mi$.
\end{proof}

\begin{theorem}[The two-row family $(2m-2,2)$]\label{thm:two}
For $m\ge2$ \emph{(}$n=2m$\emph{)},
\[
   \boxed{\,G_{(2m-2,2)}(i)=m(m-2)\,i\,},
\]
which vanishes \emph{if and only if} $m=2$, i.e.\ if and only if $\la=(2,2)$.
\end{theorem}
\begin{proof}
Put $a=2m-2$. For $a\ge4$ the removable $2$-strips of $(a,2)$ are exactly: the arm
pair $(1,a),(1,a-1)$ (horizontal, $\to(a-2,2)$, weight $1$); the second-row pair
$(2,1),(2,2)$ (horizontal, $\to(a)$, weight $1$); and the disconnected pair
$(1,a),(2,2)$ (distinct row and column since $a>2$, $\to(a-1,1)$, weight $1+i$).
These exhaust the partition-preserving deletions of two cells, the only corners being
$(1,a)$ and $(2,2)$. Hence, with $g_m:=G_{(2m-2,2)}(i)$ and using $G_{(a)}=1$
(Prop.~\ref{prop:extreme}) and $G_{(a-1,1)}=(m-2)+(m-1)i$ (Lemma~\ref{lem:21},
$M=m-1$), for $m\ge3$:
\[
   g_m=g_{m-1}+1+(1+i)\big[(m-2)+(m-1)i\big]
      =g_{m-1}+1+\big(-1+(2m-3)i\big)=g_{m-1}+(2m-3)\,i,
\]
since $(1+i)[(m-2)+(m-1)i]=(m-2)-(m-1)+[(m-1)+(m-2)]i=-1+(2m-3)i$. The base is
$g_2=G_{(2,2)}(i)=0$ (Proposition~\ref{prop:m2}). Summing,
$g_m=i\sum_{k=3}^m(2k-3)=i\big[(m-1)^2-1\big]=m(m-2)\,i$, zero iff $m=2$.
\end{proof}

\begin{corollary}\label{cor:two-row}
Among two-row shapes $(a,b)$ the value $G_{(a,b)}(i)$ vanishes only at $(2,2)$, for
$b\le2$ \emph{(}rows, $b{=}0$; near-rectangular, $b{=}1$; Theorem~\ref{thm:two},
$b{=}2$\emph{)}; the cases $b\ge3$ are verified vanishing-free for $n\le24$.
\end{corollary}

\begin{theorem}[Computational verification]\label{thm:comp}
For every $\la$ with $|\la|\le16$ (both parities; $913$ shapes through $n=16$),
$G_\la(i)=0$ if and only if $\la=(2,2)$; rechecked through $n=18$.
\end{theorem}

\section{The remaining statement and why it is hard}

By Theorem~\ref{thm:reduction} the full theorem $G_\la(i)=0\iff\la=(2,2)$ is
\emph{equivalent} to the following.

\begin{conjecture}[Full support]\label{conj:full}
For $m\ge3$, $\supp_\CC(\psi^m)=\{\la\vdash2m\}$; and for $m=2$ the support omits
exactly $(2,2)$. Equivalently, $P_\la(i)\neq0$ for all $\la\neq(2,2)$.
\end{conjecture}

The $m=2$ clause is Proposition~\ref{prop:m2}; the $m=1$ case is
$\psi=s_2+is_{11}$ (full). What remains is the $m\ge3$ clause.

\paragraph{Why positivity fails.} Writing $P_\la(i)=\Real P_\la(i)+i\,\Imag
P_\la(i)$, the real and imaginary parts are the signed sums
$\Real=\sum_{k}(-1)^{k}a_{2k}$, $\Imag=\sum_{k}(-1)^{k}a_{2k+1}$ where
$P_\la(t)=\sum_j a_j t^j$, $a_j\ge0$. Neither is sign-definite, and neither
vanishing-set is empty: e.g.\ at $m=3$, $\Real P_{(4,2)}(i)=0$ while
$\Imag P_{(2,2,1,1)}(i)=0$. The content of Conjecture~\ref{conj:full} is that the
two vanishing-sets are \emph{disjoint} (except at $(2,2)$). There is no single
Schur-positive function whose support computes $P_\la(i)$, because $\psi^m$ is
generically irreducible as a $t$-polynomial $P_\la(t)$ (for instance
$P_{(5,5)}(t)=t^5+10t^3+10t^2+15t+6$).

\paragraph{Two routes, and where each stalls.} The recursion \eqref{eq:rec}
at $t=i$ reads $P_\la(i)=\sum_{\mathrm{horiz}}P_\mu(i)+i\sum_{\mathrm{vert}}P_\mu(i)$,
a sum of Gaussian integers; the model $P_\la(i)=\sum_{\pi}i^{\,\mathrm{vert}(\pi)}$
runs over chains removing $2$-strips. At $t=-1$ the analogous signed sum
$\sum_\pi(-1)^{\mathrm{vert}(\pi)}=\chi^\la(2^m)$ is resolved by a classical
sign-reversing involution (domino tableaux), giving
$\ord_{q=-1}G_\la=\lfloor|2\text{-core}|/2\rfloor$. The tempting $\zeta_4$ analogue --
a sign-reversing $2\times2$-flip involution whose fixed set is $4$-quotient-indexed --
\emph{fails}: it is not an involution (overlapping flip windows), and the value
$\sum i^{\mathrm{vert}}$ over its non-fixed chains is a nonzero set-invariant (first
breaking at $(3,3)$, where the non-fixed chains form an unpairable $3$-cycle), so no
such involution can exist. Two routes remain.
\begin{enumerate}
\item \textbf{A $(1+i)$-adic monovariant.} It suffices that $v_{1+i}(G_\la(i))$ be
finite for $\la\neq(2,2)$. Computation confirms finiteness but shows the valuation is
\emph{irregular} (values $2,3,5,6,8,9,11,\dots$, no closed form against the
$2$-core/$4$-quotient data), so the clean $q=-1$ law does not persist at $i$. One
would need a combinatorial valuation-minimal, uniquely-placed predecessor in
\eqref{eq:rec}; I could not pin one uniformly. \emph{(}The staircase $4$-core
$(4,3,2,1)$ has $G(i)=24(1+i)^5$, a pure prime power, $v_{1+i}=11$.\emph{)}
\item \textbf{A Hermitian lower bound.} From $\psi^m\overline{\psi^m}=\Xi^m$ with
$\Xi=\psi\bar\psi=h_2^2+e_2^2=s_4+s_{31}+2s_{22}+s_{211}+s_{1^4}$ (Schur-positive,
\emph{full} support of partitions of $4$, unlike the $d=3$ $\Theta$), one has for every
$\rho\vdash4m$
\[
   0\le\inner{s_\rho,\Xi^m}=\sum_{\la,\mu\vdash2m}c^{\rho}_{\la\mu}\,G_\la(i)\,\overline{G_\mu(i)},
\]
a nonnegative Hermitian form. A $\rho$ with a \emph{unique} Littlewood--Richardson
halving $\rho=\la+\la$ would give $0<\inner{s_\rho,\Xi^m}=c^{\rho}_{\la\la}|G_\la(i)|^2$
and finish; but no such ``primitively halvable'' $\rho$ exists in general (e.g.\
$\rho=2\la$ has several halvings), so the off-diagonal terms cannot be discarded.
\end{enumerate}
The natural home for the missing positivity is the even (pair-of-boxes) Fock-space
formalism of the seed: $\psi^m=\big(\tfrac{1+i}2\big)^m(p_1^2-i\,p_2)^m$ is a genuinely
complex tau-type vector, not a positive one.

\section{Computational verification}
All identities were checked by exact computation (Gaussian integers in $\ZZ[i]$,
Schur expansions via Pieri; scripts in
\texttt{\~{}/projects/scratch/2026-06-09-d4-involution/}). Specifically: the
reduction Theorem~\ref{thm:reduction} and $G_\la(i)=P_\la(i)$ against direct
evaluation, $n\le14$ ($0$ mismatches); the integral representation \eqref{eq:gauss}
and Parseval identity \eqref{eq:parseval} for $m\le5$; the conjugation symmetry
\eqref{eq:conj}; the basic and refined parity congruences
(Prop.~\ref{prop:kostka}, Prop.~\ref{prop:refined}) for all $\la$, $n\le16$;
Propositions~\ref{prop:extreme},~\ref{prop:m2}; the closed forms
$G_{(2M-1,1)}=(M-1)+Mi$ and $G_{(2m-2,2)}=m(m-2)i$ (Lemma~\ref{lem:21},
Theorem~\ref{thm:two}) for $m\le9$; two-row shapes vanishing-free except $(2,2)$ for
$n\le24$ and hooks for $n\le22$ (Cor.~\ref{cor:two-row}); and the master scan
(Theorem~\ref{thm:comp}, $n\le16$, rechecked $n\le18$). The non-single-valuedness of
$\psi^m$ through the $2$-quotient, and the irregularity of $v_{1+i}(G_\la(i))$, were
checked for $n\le14$.

\end{document}
